\section{Challenges}
\label{sec:challenges}

Extending eBPF programmability to GPUs introduces fundamental challenges arising from (1) monolithic GPU drivers lacking designed-for-programmability interfaces, (2) fundamental mismatches between scalar eBPF semantics and GPU's SIMT execution model, and (3) complexities in coordinating consistent state across asynchronous host-device boundaries.

\subsection{C1: Safe and Efficient Driver Extensibility}

GPU drivers, such as NVIDIA's open GPU kernel modules, are large, monolithic codebases tightly integrating low-level hardware mechanisms and resource-management policies. Unlike Linux kernel subsystems explicitly designed with stable eBPF extension points (e.g., \texttt{struct\_ops}), GPU drivers today do not provide interfaces intended for safe policy programmability. Deep mechanism-policy coupling makes it challenging to introduce stable and efficient attach points for programmable policies involving critical decisions such as memory migration, prefetching, and kernel scheduling, as these are deeply embedded within vendor-specific logic like MMU page fault handlers and interrupt delivery paths. Moreover, since policy logic executes directly in privileged driver contexts with tight latency constraints and direct GPU hardware control, there are strict driver safety requirements: logic errors (e.g., incorrect lock usage or inconsistent command submission) can cause device-wide stalls or node-level failures, mandating strict bounded execution, memory safety, and driver integrity guarantees.

\subsection{C2: SIMT-Aware Device-Side Execution}

Standard eBPF is designed for scalar, single-threaded execution on CPUs, but GPU kernels operate under the fundamentally different Single Instruction Multiple Threads (SIMT) execution model, creating significant programmability mismatches. Naively mapping scalar eBPF handlers onto GPU threads introduces severe warp divergence overhead, significantly degrading GPU throughput and occupancy. For instance, executing scalar policy logic independently on all 32 threads in a warp increases instruction divergence by up to 32$\times$, drastically reducing kernel performance. Furthermore, GPU architectures feature complex memory hierarchies (per-thread registers, shared memory, L2 caches, global memory) and weak memory consistency semantics (\texttt{\_\_threadfence\_system}), posing memory safety verification challenges fundamentally unaddressed by existing CPU-side eBPF tools under thousands of concurrent warp accesses. Lastly, GPU kernels execute within flat address spaces without isolation or kernel-level runtime recovery mechanisms; thus, policy errors like infinite loops or invalid memory accesses cause unrecoverable hardware stalls, making bounded execution and memory safety guarantees essential to enforce statically at compile time.

\subsection{C3: Unified Cross-Layer State and Coordination}

Many GPU policies require coordinated actions spanning CPU host drivers and GPU device execution; for example, GPU-side instrumentation capturing warp-level memory accesses informs host-side memory placement and scheduling decisions. Ensuring efficient, safe, and consistent cross-layer coordination poses unique technical difficulties. First, there is a significant cross-boundary memory hierarchy mismatch: placing shared maps naively in host memory introduces prohibitive latency (hundreds of microseconds per access from GPU warps), drastically impacting GPU utilization, whereas placing all shared state within GPU global memory risks limited capacity, bandwidth bottlenecks, and cache thrashing, requiring automatic balancing of latency, capacity, and access locality. Additionally, host driver operations and GPU execution proceed asynchronously at vastly different granularities, complicating state consistency. Without explicitly defined consistency abstractions, host-side policy decisions (e.g., page eviction or kernel scheduling) risk relying on outdated or inconsistent GPU-side signals (e.g., warp-level memory access statistics), causing incorrect memory migration or scheduling decisions and significantly reducing system effectiveness.

\section{\sys Design}

\subsection{Design Principles}

\sys follows three key design principles to address the limitations of existing GPU stacks:

\paragraph{Principle 1: Decouple stable mechanisms from dynamic programmable policies.}
\sys factors the GPU stack into a small set of stable, privileged mechanisms and a separate, dynamically programmable policy layer. The \sysdriver exposes 2 level of policy mechanisms: driver is for transparent global, fine-grained policy that must be ensured across process and tenant, while device policies include per-application level fine-grained control and performance hints. All decisions live in eBPF policies attached to these hooks. We mandate driver modification because purely user-space or sidecar approaches lack the privileges to interpose on critical hardware paths (e.g., MMU faults, context switching) and cannot enforce global, cross-tenant invariants.

\paragraph{Principle 2: Unified CPU--GPU policy interface and control plane.}
All \sys policies, whether running in the OS kernel or inside GPU kernels, are written in the same restricted eBPF IR and use the same abstractions. A single control plane loads, verifies, and attaches eBPF programs, configures shared maps, and coordinates their lifecycle. This sacrifices some expressiveness compared to arbitrary CUDA code (e.g., direct kernel manipulations), but eliminates cross-layer fragmentation and makes it straightforward to split a policy across host and device.

\paragraph{Principle 3: Safety and efficiency through verified bounded execution and interfaces}
\sys only executes policy code that is expressed in a restricted eBPF subset and accepted by the in-kernel verifier, each type with certain kfunc. This design intentionally trades policy expressiveness for strong guarantees and efficiency: policy handlers cannot hang the GPU, corrupt driver or application memory, introduce unbounded latency on critical paths or use too much resource that hurt the application.

\subsection{High-Level Architecture}
\label{sec:arch}

\begin{figure}
\centering
\fbox{\parbox{0.8\columnwidth}{\centering [Figure: \sys Design]}}
\vspace{-.5ex}
\caption{\sys high-level design: unified eBPF runtime across CPU and GPU.}
\label{fig:system-design}
\end{figure}

Figure~\ref{fig:system-design} shows the high-level architecture of \sys, including the CPU--GPU boundary, \sysdriver extensions, and GPU kernel-level runtime integration points. At its core, \sys provides a unified eBPF runtime that manages verification and just-in-time (JIT) compilation in the operating system kernel or the userspace loader.

Policy authors write \sys programs using standard eBPF tooling (e.g., \texttt{clang} with \texttt{libbpf}, \texttt{bpftool}, or \texttt{bpftrace}). Programs are compiled to eBPF bytecode, loaded into the kernel, and verified by our extended kernel verifier before execution. \sys loader places eBPF maps in memory regions visible to both the OS kernel and GPU kernels.

Once verified, the same bytecode can be installed in two ways:

\begin{itemize}[leftmargin=*]
  \item \emph{Host side:} \sys attaches eBPF programs to hooks in \sysdriver at critical policy decision points in the GPU stack (e.g., memory prefetching, migration, and kernel scheduling), or use other eBPF CPU hooks. These hooks follow a \texttt{struct\_ops}-style callback interface: when the driver reaches a hook, it invokes the attached eBPF handler with an event-specific context structure. The handler returns a compact decision (e.g., an action code or priority) that the driver enforces or using safe kfunc calls into the driver code.
  \item \emph{Device side:} for execution inside GPU kernels, \sys JIT runtime inside the loader translates the verified eBPF bytecode to device bytecode and injects it into the target GPU context. JIT runtime use a pass like pipeline abstraction to allow user do instructions level instrucmentation to the GPU device bytecode, define new attach points, insert the function calls before and after any instruction, and custom define their optimizations. Lightweight trampolines inserted at kernel entry or other selected sites save minimal live state, call into the translated eBPF snippet, and then restore state to resume the original control flow. Device-side eBPF programs update the same maps as their host-side counterparts.
\end{itemize}

A GPU-side runtime enforces the safety guarantees described in the Safety Model (\S\ref{sec:prog-model}): each injected program operates in a protected domain that cannot interfere with application kernels, while still enabling adaptive policy decisions on critical GPU execution paths.

Shared eBPF maps provide efficient state coordination between \sysdriver and eBPF policy logic on the host, on the device, and between host and GPU policies across the CPU--GPU boundary. This design yields real-time, low-overhead cross-layer observability and control while retaining the safety guarantees of the eBPF framework. We focus on multi-tenant policy control on shared GPUs, not on providing virtual GPU abstractions or hardware partitioning. Mechanisms such as MIG or SR-IOV are orthogonal and can, in principle, be deployed underneath \sys; \sys policies only manipulate well-defined \sys-managed state and cannot corrupt kernel or application memory outside that state.


\subsection{\sys Programming Model}
\label{sec:prog-model}

A \sys policy is a distributed application consisting of eBPF code that runs in
privileged contexts on the CPU and GPU device together, with a user-space control-plane
component. All pieces program against a unified abstraction: a restricted eBPF IR,
a small set of \sys-specific attach points and helper functions different for each attach types, and shared eBPF maps and metadata.

\subsubsection{Policy Code (CPU \& GPU)}

\sys introduces two primary eBPF program types to categorize policy logic:
a driver-level type for GPU policy hooks (encompassing both \texttt{BPF\_PROG\_TYPE\_GPU\_MEM} for memory events and \texttt{BPF\_PROG\_TYPE\_GPU\_SCHED} for scheduling) and a device-level type (\texttt{BPF\_PROG\_TYPE\_GPU\_DEV}) for GPU-kernel instrumentation. Each program type is associated with a specific family of hook contexts that are rigorously defined and BTF-annotated (BPF Type Format) to be understood by the kernel verifier. Each program defines multiple entry points corresponding to attach points on either the host or GPU device.

On the host, verified eBPF bytecode executes directly in privileged kernel contexts provided by \sysdriver, allowing policies to respond to critical memory events (such as device memory misses, eviction candidate selection) and kernel scheduling decisions (kernel submission, admission, preemption, completion). On the GPU, the same verified bytecode is translated to device instructions (e.g., PTX or SPIR-V) and executed through lightweight trampolines injected at kernel entry points, selected memory operations, and synchronization or phase boundaries. GPU side eBPF can only use helpers that implemented as part of the GPU trampoline. In both host and GPU contexts, policy handlers receive structured event contexts containing metadata (e.g., virtual addresses, region IDs, kernel IDs, warp identifiers), and return compact policy decisions enforced directly by underlying GPU mechanisms.

\subsubsection{Shared State \& Metadata}

Shared eBPF maps store policy-relevant state including counters, histograms, per-region hotness metrics, per-tenant scheduling priorities, and runtime configuration parameters. This enables device-side probes, such as GPU memory access instrumentation or warp progress signals, to update policy metadata that host-side memory or scheduling hooks later consult. \sys transparently propagates and translates core metadata such as virtual memory addresses, kernel identifiers, timestamps across CPU and GPU, presenting a logically unified identifier space. This design allows policy developers to implement sophisticated cross-layer logic (e.g., GPU-side access tracking coupled with host-side eviction or scheduling) without manually handling address translation or cross-boundary synchronization complexities.

\subsubsection{Control-Plane Application and Lifecycle}

A \sys policy also includes a dedicated user-space control-plane component responsible for lifecycle management, dynamic reconfiguration, and policy introspection. This component instantiates and manages eBPF maps, and attaches policy sections to defined host and GPU device hook points exposed by \sysdriver and GPU-side runtime. It enables dynamic reconfiguration of deployed policies at runtime by updating parameters (e.g., thresholds, sampling rates, eviction aggressiveness, power caps) stored in shared maps, and facilitates swapping or updating of attached policy logic without restarting applications or kernels. Furthermore, the control-plane application continuously monitors policy operation by reading metrics and state from maps, such as memory hotness counters, per-tenant scheduling metrics, and queueing latencies—and exports this information to external observability frameworks or adaptation controllers. Thus, from a developer's perspective, a complete \sys policy application comprises policy logic for both CPU and GPU execution contexts along with a complementary control-plane component orchestrating policy behavior and evolution over time.

\subsection{Policy Interfaces}
\label{sec:design-interfaces}

\subsubsection{Memory Interface}
\label{sec:design-mem}

% Modern GPU memory subsystems move data across device memory, host memory, and
% intermediate tiers using fixed, internally implemented heuristics, often tied to
% page faults or simple recency-based policies. These mechanisms neither expose their
% decisions nor exploit workload-specific structure, and they interact poorly with
% user-space caches that increasingly manage model weights, KV-cache segments, and
% expert parameters.
\sys treats memory placement as a programmable region
cache: it exposes a small set of region-level events for \emph{admission},
\emph{access}, and \emph{removal}, together with a prefetch channel, and lets host
and device cooperate on a unified policy.

On the host driver, these events are realized as a small \texttt{struct\_ops}-style policy
table with typed contexts:

\begin{minted}[highlightlines=false,fontsize=\small]{c}
// Host-side policy callbacks for GPU memory management.
struct gdrv_mem_ops {
  // Called when a region is created or becomes eligible for device placement.
  int (*region_add)(gdrv_mem_add_ctx *ctx);
  // Called when memory subsystem observes region use; decide cache hit,
  int (*region_access)(gdrv_mem_access_ctx *ctx);
  // Called when a region is about to leave the device-resident working set.
  int (*region_remove)(gdrv_mem_remove_ctx *ctx);
  // Called at safe points (between kernel launches, page faults) for prefetch.
  int (*prefetch)(struct gdrv_mem_prefetch_ctx *ctx);
};
// Host-side kfunc, operate on a block (virtual address unit) across GPU and cpu
bool gdrv_mem_list_insert_head(struct gdrv_mem_list* head, struct gmem_block* r);
bool gdrv_mem_list_insert_tail(struct gdrv_mem_list* tail, struct gmem_block* r);

// device-side memory callback interface
struct gdev_mem_ops {
    void (*memtrace)(...);
    void (*memsync)(...);
    ...
};
// device side helpers
// hint to ask for kernel prefetch and trigger the prefetch callback in kernel
int gdev_mem_prefetch_hint(struct gmem_region* r);
// hint to pin a memory on CPU or on GPU
int gdev_mem_pin_hint(struct gmem_region* r, int device_id);
\end{minted}

\paragraph{Host-side region events and removal.}
On the host, \sys models memory as a set of logical regions and exposes attach points for three
canonical cache events. To ensure safety, e.g. no region is leak or always pin, the regions is organized as a list: it's always added from the end of the list and removed from the head of the list. The eBPF program can reorder the list, preserve state and addtional data with list nodes to express complex algorithms.

Together, these events realize a
general region-cache interface: policies implemented as eBPF handlers maintain their
own state in shared maps (for example, LRU lists, LFU counters, or tenant budgets),
and respond to Add/Access/Remove and prefetch callbacks with compact decisions
such as \emph{admit}, \emph{bypass}, \emph{pin}, \emph{evict}, or ordered lists
of regions to migrate. The underlying memory subsystem enforces these decisions
and maintains correct mappings, but delegates replacement and prefetch policy to
\sys.

\paragraph{device-side memory events.}

On the GPU, \sys exposes a single device-side memory event interface that kernels
use both to report observed accesses and to express future-use hints. Device-side
handlers receive a compact context struct (not shownp') with fields such as a logical
region identifier, kernel, warp and SM identifiers, an event type (\emph{ACCESS}
or \emph{HINT}), and optional hint flags.

At selected memory operation sites, \sys injects lightweight hooks that operate at
warp granularity and invoke a device-side \sys handler with an \emph{ACCESS} event:
the handler updates shared maps with per-region statistics such as access counts,
approximate working-set membership, and simple locality signatures.\

The same device-side interface also supports \emph{HINT} events, which kernels
invoke explicitly when they have domain knowledge about future region usage. For
example, a mixture-of-experts kernel can compute the identifiers of experts it
will access for the next batch and call a \sys helper with HINT events that mark
the corresponding regions as "will-use soon" with an optional priority or deadline;
a KV-cache kernel can emit HINT events that mark segments as cold and suitable for
demotion. Both ACCESS and HINT events are recorded in shared maps or small control
queues that host-side \sys handlers consult at Access, Remove, and prefetch hooks.
Device-side code never manipulates page tables or performs migration directly; it
implements the sensing and hinting side of a closed-loop policy, while host-side
\sys handlers implement the placement, prefetch, and eviction decisions over the
GPU memory hierarchy.


\subsubsection{Scheduling Interface}
\label{sec:design-sched}

GPU schedulers today typically implement fixed, coarse-grained policies such as FIFO
queues with limited priority control, which are insufficient in multi-tenant settings
with mixed latency-sensitive and batch workloads or tight power constraints. In such
environments, operators need policies that can shape both \emph{which kernels reach
the GPU} at the granularity of kernel, and \emph{how work is consumed inside those kernels} at the granularity of thread blocks or tiles. \sys exposes scheduling interfaces on both host and device
so that policies can implement fine-grained priority, deadline-aware, and power-aware
scheduling while respecting the GPU SIMT execution model.

On the host, these lifecycle events are exported as a small struct-ops table of
scheduling callbacks:

\begin{minted}[highlightlines=false,fontsize=\small]{c}
// Host-side policy callbacks for GPU kernel scheduling.
struct gdrv_sched_ops {
  // Called when a kernel is enqueued by a runtime.
  // Return 0 to accept, <0 to reject or defer.
  int  (*submit)(struct gsched_submit_ctx *ctx);
  // Called when the scheduler is about to dispatch a ready kernel.
  // Returns the ID of the kernel to admit, or -1 to defer.
  int  (*admit)(struct gsched_admit_ctx *ctx);
  // Notification that a kernel has completed. Must not fail.
  void (*complete)(struct gsched_complete_ctx *ctx);
};
// Host side kfunc
bool gdrv_sched_premmpt(struct gched_preempt_ctx *ctx);
// Device-side block scheduling callbacks.
struct gdev_sched_ops {
  // Called when a worker block starts a new work unit.
  void (*entry)(struct gdev_block_ctx *ctx);
  // Called when a worker block finishes a work unit.
  void (*exit)(struct gdev_block_ctx *ctx);
  // Called when a worker block is idle; returns whether
  // it should try to steal work, as fine-grain block level scheduler
  bool (*should_try_steal)(struct gdev_block_ctx *ctx);
};
\end{minted}

\paragraph{Host-side scheduling hooks.}
On the host, \sys extends the GPU kernel lifecycle with attach points that allow
policies to control admission, ordering, and preemption of kernels:

Attached eBPF schedulers read these contexts and shared maps (for example, per-tenant
latency and throughput statistics, outstanding work across queues) and return compact
decisions such as which kernel to run next, whether to throttle a tenant, or whether
to allow preemption. The host-side interface thus controls which kernels occupy GPU
contexts and how they are multiplexed in multi-tenant scenarios.
\paragraph{Device-side thread-block scheduler and signals.}

Inside the GPU, \sys exposes a complementary interface that operates at the level of
thread blocks or tiles, aligning with the SIMT execution model. Each kernel that uses \sys's internal scheduler is modeled as a pool of \emph{work units} (e.g., tiles or logical block indices) consumed by \emph{worker blocks}. Device-side policies observe and steer this internal scheduler via a \texttt{gdev\_block\_ctx} that includes the kernel identifier, tenant identifier, and a logical work-unit or block identifier.

Together, the host-side and device-side interfaces allow \sys policies to coordinate decisions across the full scheduling pipeline: from kernel admission and preemption on the host to fine-grained block-level work distribution inside GPU kernels.

\paragraph{Example: priority-aware block scheduler.}
As a concrete illustration, consider a policy that gives low latency to an interactive tenant while preserving throughput for batch jobs. The host keeps a per-tenant record of outstanding work and virtual time in a BPF map. \texttt{submit} tags kernels with a tenant ID and initial priority. \texttt{admit} prefers kernels from the interactive tenant when its queue is non-empty and boosts its priority when its p99 latency (tracked in another map) exceeds a target. On the device, each worker block's \texttt{entry} handler increments a per-tenant ``active blocks'' counter, and \texttt{exit} decrements it and records the work-unit completion time. When \texttt{should\_try\_steal} runs on an idle block, it uses these maps to decide whether to steal a work unit from the interactive tenant or from the batch pool. In \S\ref{sec:eval} we show that this policy reduces p99 latency by 2--5$\times$ with 8--12\% throughput loss for batch work.

\subsection{\sys GPU device SIMT specific verification and optimizations}
\label{sec:safety-model}

\paragraph{Safety Model and TCB.}
The core challenge in programmable GPU policy is to run operator-supplied logic in privileged, latency-critical GPU driver and device paths without process-like isolation: even simple bugs can corrupt mappings, stall the device, or inflate tail latency, which even cause cluster level fault in critical training process. \sys runs operator-supplied policy code in these privileged paths. We assume the entity deploying \sys policies (e.g., a cluster operator) is trusted but fallible: policies may contain logic bugs or misconfigurations but are not malicious. Malicious kernel modules, vendor drivers, firmware, or hardware are out of scope.

Under this model, the \sys TCB consists of: (1) the in-kernel eBPF verifier, (2) the \sys kernel module, (3) the device-code backend, and (4) the GPU driver/firmware. Policy authors only reason about eBPF and the \sys API; they never write PTX/SASS directly. All policies are expressed in a restricted eBPF subset and must pass the in-kernel verifier, which enforces bounded loops, stack and map safety, and restricts helpers and kfuncs to the \sys-specific set.

\paragraph{SMIT aware Verification}

The driver side is simple, we simply use the kernel's eBPF verifier, checks that each program uses only the \sys-specific kfuncs and hook contexts declared via BTF annotations. The kernel verifier enforce:
1. memory safety
2. type Safety
3. resource safety
4. bound and no hardware exceptions

Besides this,

So, we explicitly extend kernel eBPF verifier through kernel module to:
1. use conservative work budgets for execution instructions, memory usage, and memory access count per instructions, to avoid memory bandwidth overload. Define per-program-type budgets that reflect how often hooks fire and what overhead you can tolerate, e.g. \texttt{insn\_bound} $\le$ \texttt{B\_max\_prog\_type}, \texttt{map\_ops} $\le$ \texttt{M\_max\_prog\_type}, branch $<$ \texttt{M\_max\_prog\_type}
2. avoid sync and atomic, global barriers or device-wide synchronization, ban the eBPF instructions
3. enforce the new and GPU specific helpers
3. what is more ?

There is no runtime sandbox; if the program gets past the verifier, it runs natively on GPU; if it doesn't, it never gets loaded.

\paragraph{SIMT aware optimization}

Naively executing scalar control logic across all GPU threads introduces prohibitive divergence and overhead. The \sys JIT implement a optimization pass pipeline, allow:
1. \sys aggregates scalar policy logic execution at warp granularity, executing decisions on a designated lane while synchronizing minimal input from other lanes.
2. We find a pattern that most maps does not need golabl access: mutl level map design and auto placement: Frequent global state updates are optimized through warp-level and SM-level local maps dramatically reducing contention and costly global memory operations. logically shard, but not actually shared. user can mark a map and move it from cpu/gpu memory explictly. we model the eBPF map access tp different place as distance, similar to the ideqqa of numa node, and by static analysis the access of eBPF instructions, we can calc the best place to put the maps to.
4. so, unlike kernel ebpf, in gpu functions and maps after verification are aggressively inline to avoid call cost, and we restrict the gpu map size to fit into the register..
5. since we can inline the maps and helpers, we can eliminates branch cost as much as possible by remove the if after
6. allow user to config sampling to reduce overhead. (Do we need this?)
Collectively, these optimizations enable efficient, scalable, and minimally invasive GPU-side policy execution suitable for production deployment.
